\chapter{Adding a prophet: building PUCT one failed attempt at a time}
\label{ch:puct}

\begin{objectives}
\begin{itemize}
\item Say what a policy prior is and what it is not.
\item Build the PUCT selection rule through five attempts, and be able to say which failure
  each surviving symbol repairs.
\item Prove the equilibrium property: when search learns nothing, visits become proportional
  to the prior --- and check it against real output.
\item Compute $U(s,a)$ exactly as Leela computes it, including the growth of $c_{\text{puct}}$
  and the first-play-urgency value of an unvisited move.
\item Verify your arithmetic against the engine's own \texttt{S} column to five decimals.
\end{itemize}
\end{objectives}

\section{The prophet}

Chapter~\ref{ch:uct} ended stuck: UCT must taste every move before it can prefer any, and at
$b\approx31$ that is unaffordable. What we want is what a strong player has --- a glance that
says \emph{these two moves, probably; that one, conceivably; the other twenty-eight, no}.

Assume we have one. Formally, a \textbf{policy prior} is a probability distribution over the
legal moves of a position:
\[
  P(a \given s) \ge 0, \qquad \sum_{a \in \mathcal{A}(s)} P(a\given s) = 1 .
\]
Where it comes from is Part~\textsc{iv}'s problem; here it is an oracle we may consult for
free. What it means is worth being exact about, because the whole design depends on it:

\begin{keyidea}
$P(a\given s)$ is \textbf{not} an estimate of how good the move is. It is an estimate of
\emph{how likely this move is to be the one finally played} --- literally, what fraction of
the time a strong player (in Leela's case, its own trained self) would choose it. It is a
\textbf{claim about attention}, not about value. A move can be excellent and have low $P$ if
it is the sort of move that is usually wrong.
\end{keyidea}

That last sentence is not a technicality. It is the mechanism behind everything your project
calls a "hidden gem", and here it is in real numbers before we start:

\begin{realdata}[title={Real engine output: priors in the Morphy position, before 16.Qb8+}]
Position after 15\ldots Nxd7 in Morphy--Brunswick/Isouard, Paris 1858. 46 legal moves. Top
priors from the network in \texttt{engine/791556.pb.gz}, no search:
\begin{center}
\begin{tabular}{@{}lrl@{}}
\toprule
Move & $P(a\given s)$ & \\
\midrule
Qb7   & 31.93\% & the network's instinct \\
Qb5   &  9.57\% & \\
Qc3   &  8.33\% & \\
Qxe6+ &  7.92\% & \\
Rxd7  &  7.81\% & \\
Qa4   &  7.65\% & \\
\ \ \vdots & \vdots & \\
\textbf{Qb8+} & \textbf{1.60\%} & \textbf{Stockfish (depth 30): mate in 2} \\
\bottomrule
\end{tabular}
\end{center}
The move that forces mate in two gets $1.60\%$ of the network's attention. It is not ranked
first, second or third. Any design that decides in advance which moves are worth looking at
must survive this position.
\end{realdata}

\section{Five attempts}

We want a selection rule of the shape $\argmax_a [\,Q(s,a) + U(s,a)\,]$ where $U$ now uses
$P$. Chapter~\ref{ch:uct} set three requirements: (R1) most moves may get \emph{zero} visits;
(R2) any move must still be reachable later; (R3) with enough visits, measurement must be
able to overrule the prior.

\begin{attempt}{1}{filter the move list --- search only the top $k$ priors}
Simple, fast, and used by real programs in other domains. Take $k = 3$ in the Morphy
position: we search Qb7, Qb5, Qc3, and \textbf{mate in two is not in the list}. Not at 800
nodes, not at eight billion.

Fails R2, catastrophically and silently. A hard cutoff converts the prior's mistakes into
the engine's permanent blind spots --- and the moves a policy under-rates are exactly the
surprising ones, which are exactly the ones a training tool exists to show you.
\end{attempt}

\begin{attempt}{2}{add the prior to the value --- $Q(s,a) + c\,P(a\given s)$}
Now every move is reachable in principle. But the prior's contribution never decays, so it
is not a \emph{bonus for uncertainty}, it is a permanent thumb on the scale.

With $c = 1$ in the Morphy position: Qb7 gets a standing $+0.319$ and Qb8+ a standing
$+0.016$. Suppose the search does 100{,}000 visits and establishes the true values --- Qb8+
wins by force, $Q = 1.0$; Qb7 is merely much better, $Q \approx 0.66$. Scores: Qb8+ gives
$1.016$, Qb7 gives $0.979$. Here it happens to work, because the truth gap ($0.34$) exceeds
the prior gap ($0.30$). Make the position one where the best move wins by a small margin and
the prior wins permanently.

Fails R3: no amount of evidence can wash the prior out. And it fails R1 too, in the other
direction --- because the term is a constant, the search has no reason to \emph{stop}
returning to a well-measured move.
\end{attempt}

\begin{attempt}{3}{multiply the prior into UCB1's bonus ---
  $Q + c\,P\sqrt{\dfrac{\ln N}{n}}$}
Now we are close to something sensible: the bonus decays with visits (good), grows slowly
with total search (good), and is scaled by how promising the move looked (new and good).

One fatal defect remains, and it is the same one that sank UCT: at $n = 0$ the term
$1/\sqrt{n}$ is \textbf{infinite}. Every legal move still gets a mandatory first visit,
whatever its prior. We are back to spending 31 visits before the second look at anything.

Fails R1. The fix must make the bonus \emph{finite} at zero visits.
\end{attempt}

\begin{attempt}{4}{make it finite at zero --- $Q + c\,P\dfrac{\sqrt{\ln N}}{1+n}$}
Replacing $1/\sqrt{n}$ with $1/(1+n)$ does exactly what is needed: at $n=0$ the bonus is
$c\,P\sqrt{\ln N}$, which is \emph{small for a move with a small prior}. A move with
$P = 0.4\%$ starts with a small claim on the budget, and the search is free to ignore it.
R1 satisfied.

But now look at the cost of that change. The denominator went from $\sqrt{n}$ to $n$, so the
bonus decays much faster than before. Meanwhile the numerator still grows only like
$\sqrt{\ln N}$ --- glacially: from $N = 100$ to $N = 10{,}000$ it grows by a factor of
$\sqrt{\ln 10^4/\ln 10^2} = \sqrt{2} = 1.41$. Fast decay plus glacial growth means a
neglected move's priority stays flat essentially forever. In the Morphy position, Qb8+'s
bonus would creep up by $41\%$ over a hundredfold increase in search --- while it needs to
grow by a factor of about ten to become competitive.

Fails R2 in practice, if not quite in theory. \textbf{The numerator has to grow faster.}
\end{attempt}

\begin{attempt}{5}{let the numerator grow like $\sqrt{N}$}
\[
  U(s,a) = c_{\text{puct}}\;P(a\given s)\;\frac{\sqrt{N(s)}}{1 + N(s,a)} .
\]
Now, from $N = 100$ to $N = 10{,}000$ the numerator grows tenfold, not by $41\%$. A move
passed over early is not passed over forever: its claim strengthens steadily as the search
pours visits elsewhere. R2 satisfied --- and we will watch it happen, in real output, in §6.5.
\end{attempt}

\begin{finalform}[PUCT --- the rule Leela actually uses]
\begin{equation}
  \label{eq:puct}
  \boxed{\;
  a^{*} \;=\; \argmax_{a}\left[\;
  \underbrace{Q(s,a)}_{\text{what I measured}}
  \;+\;
  \underbrace{c_{\text{puct}}\;P(a\given s)\;\frac{\sqrt{N(s)}}{1+N(s,a)}}_{U(s,a):\ \text{what I haven't checked}}
  \;\right]\;}
\end{equation}
with $N(s) = \sum_b N(s,b)$ the total visits already spent at this node.
\end{finalform}

Read the three moving parts against the three failures they repair:

\begin{center}
\begin{tabular}{@{}lll@{}}
\toprule
Symbol & Says & Repairs \\
\midrule
$P(a\given s)$ & "this much of my attention is owed to this move" & Attempt 1's blindness \\
$\dfrac{1}{1+N(s,a)}$ & "\ldots divided by how much I have already looked" & Attempt 3's mandatory first visit \\
$\sqrt{N(s)}$ & "\ldots multiplied by how much thinking has happened here" & Attempt 4's frozen-out move \\
\bottomrule
\end{tabular}
\end{center}

\begin{keyidea}
$U$ is a \textbf{claim on the budget}. Read $P\sqrt{N}/(1+n)$ as: \emph{"I was promised
roughly a $P$ share of the attention here; $N$ units of attention have been spent; I have
had $n$. How overdue am I?"} The whole formula is one sentence: \textbf{take the move with
the best combination of proven performance and unpaid attention.}
\end{keyidea}

\section{The equilibrium property: visits become the policy}

Here is a small theorem that makes PUCT much easier to think about --- and, unusually for
this subject, it is two lines.

\begin{established}
\textbf{Claim.} Suppose the search learns nothing: every move's $Q$ is identical (call it
$q$). Then PUCT allocates visits in proportion to the prior, $N(s,a) \propto P(a\given s)$.

\smallskip
\textbf{Proof.} Selection always picks the maximum of $Q + U$. If all $Q$ are equal to $q$,
the maximum of $q + U$ is the maximum of $U$. A rule that repeatedly picks the largest $U$
and thereby decreases it drives all the $U$'s towards equality. So at equilibrium, for any
two moves $a$ and $b$,
\[
  c_{\text{puct}}P_a\frac{\sqrt{N}}{1+n_a} \;=\; c_{\text{puct}}P_b\frac{\sqrt{N}}{1+n_b}
  \quad\Longrightarrow\quad \frac{P_a}{1+n_a} = \frac{P_b}{1+n_b}
  \quad\Longrightarrow\quad 1+n_a \;\propto\; P_a . \qquad\blacksquare
\]
\end{established}

\begin{keyidea}
\textbf{The visit distribution is the prior, reshaped by what search discovered.} If search
finds nothing, visits reproduce the policy exactly. Every departure of the visit distribution
from the policy distribution is \emph{information the search added}. That single sentence is
the foundation of Chapter~\ref{ch:readingtree}, of your project's "optical trap" and
"hidden gem" categories, and of how Leela trains itself (Chapter~\ref{ch:heads}).
\end{keyidea}

We can test the claim on real output. Take a position where the search genuinely learns
nothing --- a dead-drawn king-and-pawn ending in which every move evaluates to exactly $0.000$:

\begin{realdata}[title={Real engine output: 800 nodes on a position where all moves draw}]
White Ke4, pawn e5, Black Ke6 (Black holds the opposition). Stockfish at depth 30 confirms
$0.00$ for every move. Leela's search after 800 nodes:
\begin{center}
\begin{tabular}{@{}lrrrr@{}}
\toprule
Move & $P$ & $N$ & share of visits & $Q$ \\
\midrule
Kf4 & 19.82\% & 172 & 21.5\% & $0.0059$ \\
Kd4 & 19.74\% & 171 & 21.4\% & $0.0059$ \\
Kd3 & 20.45\% & 155 & 19.4\% & $0.0000$ \\
Kf3 & 20.19\% & 152 & 19.0\% & $0.0000$ \\
Ke3 & 19.79\% & 148 & 18.5\% & $0.0000$ \\
\bottomrule
\end{tabular}
\end{center}
Priors are nearly uniform ($\approx 20\%$ each) and so are the visits ($\approx 20\%$ each),
as the theorem predicts. The residual tilt towards Kf4 and Kd4 is not noise: those two moves
have $Q = 0.0059$ against $0.0000$ for the rest --- a difference in the fifth decimal place of
a drawn position, and the search still spent $3\%$ more of its budget on them. The mechanism
is exactly as advertised, down to a rounding error's worth of "advantage".
\end{realdata}

And now the same test where the search \emph{does} learn something:

\begin{realdata}[title={Real engine output: 1600 nodes on the initial position}]
\begin{center}
\begin{tabular}{@{}lrrrrr@{}}
\toprule
Move & $P$ & $N$ & visit share & $Q$ & share $-$ prior \\
\midrule
e4  & 22.22\% & 532 & 35.6\% & $0.1058$ & $+13.4$ pts \\
d4  & 18.41\% & 451 & 30.2\% & $0.1062$ & $+11.8$ pts \\
Nf3 & 14.22\% & 230 & 15.4\% & $0.0995$ & $+1.2$ pts \\
c4  &  9.19\% & 100 &  6.7\% & $0.0849$ & $-2.5$ pts \\
g3  &  7.26\% &  64 &  4.3\% & $0.0733$ & $-3.0$ pts \\
e3  &  4.43\% &  28 &  1.9\% & $0.0483$ & $-2.5$ pts \\
\bottomrule
\end{tabular}
\end{center}
Here the $Q$'s differ, so the visit distribution is a \textbf{sharpened} version of the
prior: the top two moves take a larger share than their priors promised, and everything below
Nf3 takes less. Search has concentrated attention on what measurement supports. That
sharpening is the search's entire contribution, and Chapter~\ref{ch:heads} shows it is also
precisely the training signal that improves the network.
\end{realdata}

\section{The two remaining pieces of real arithmetic}

Equation~\eqref{eq:puct} is the formula in every paper. The engine on your disk uses two
refinements, and you need both to reproduce its numbers exactly --- which we are about to do.

\subsection{$c_{\text{puct}}$ grows with the search}

Leela does not use a constant. It uses
\begin{equation}
  \label{eq:cpuctgrow}
  c_{\text{puct}}(N) \;=\; \texttt{CPuct} \;+\; \texttt{CPuctFactor}\cdot
  \ln\!\left(\frac{N + \texttt{CPuctBase} + 1}{\texttt{CPuctBase}}\right),
\end{equation}
with the defaults on your engine being $\texttt{CPuct} = 1.745$,
$\texttt{CPuctBase} = 38739$, $\texttt{CPuctFactor} = 3.894$ (read off
\texttt{lc0.exe}'s own \texttt{uci} option list).

The logarithm means that for small searches the growth is negligible --- at $N = 182$,
$c_{\text{puct}} = 1.745 + 3.894\ln(1.00470) = 1.745 + 0.018 = 1.763$ --- while at
$N = 10^6$ it reaches about $14$. The design intent is legible: a small search should trust
the network, and a huge search should be increasingly willing to look at what the network
dismissed. (Compare your answer to Exercise~\ref{ex:ch4:tuningc}.)

\subsection{First-play urgency: what $Q$ does an unvisited move have?}

Equation~\eqref{eq:puct} needs $Q(s,a)$ for every candidate --- including moves never visited,
which have no average at all. What value should they be given?

\begin{itemize}
\item \emph{Optimistic} ($Q = +1$): every move looks winning until checked, so everything gets
  visited once. That is UCT again.
\item \emph{Neutral} ($Q = 0$): reasonable, but wrong in won and lost positions --- in a
  position where you are winning by a queen, "0" is a pessimistic slander of every unexamined
  move; in a lost position it is wild optimism.
\item \emph{Inherit the parent} ($Q = Q(s)$): assume a move is as good as the position it
  comes from --- sensible, but then unvisited moves tie with measured ones and get pulled in
  too eagerly.
\end{itemize}

Leela's default, \textbf{FPU reduction}, takes the third option and subtracts a penalty that
grows with how much of the prior has already been explored:
\begin{equation}
  \label{eq:fpu}
  Q_{\text{FPU}}(s) \;=\; Q(s) \;-\; \texttt{FpuValue}\cdot\sqrt{\textstyle\sum_{b \,:\, N(s,b)>0} P(b\given s)},
  \qquad \texttt{FpuValue} = 0.33 .
\end{equation}
The reasoning: if you have already examined moves covering $80\%$ of the prior mass and none
of them beat what you have, the remaining $20\%$ is probably not hiding a better move ---
so discount it. Early in a node's life the penalty is near zero; late in its life it is
substantial.

\begin{worked}{deriving FPU from real output --- a piece of detective work}
Leela's verbose output prints \texttt{Q: 0.0} for an unvisited move (there is no average to
print), but its \texttt{S} column is the real PUCT score, and $\texttt{S} = Q_{\text{FPU}} +
U$. So we can recover the hidden value the engine actually used.

In the Morphy position at $N(s) = 182$, the line for Qb8+ reads
$\texttt{U} = 0.37909$, $\texttt{S} = 0.69780$. Therefore
\[
  Q_{\text{FPU}} = \texttt{S} - \texttt{U} = 0.69780 - 0.37909 = 0.31871 .
\]
Check it against Equation~\eqref{eq:fpu}. The root's own value is $Q(s) = 0.6117$, and the
engine helpfully prints the visited prior mass on its root line: $\texttt{P: 79.96\%}$. So
\[
  Q(s) - 0.33\sqrt{0.7996} \;=\; 0.6117 - 0.33(0.8942) \;=\; 0.6117 - 0.2951 \;=\; 0.3166 .
\]
Against the recovered $0.31871$ --- agreement to three decimals, the residue being the
rounding of the printed percentages. \textbf{Equation~\eqref{eq:fpu} is confirmed on real
output.}
\end{worked}

\section{Watching R2 work: the queen sacrifice comes back}

We can now close the loop on the requirement that killed three of our five attempts. Recall
Qb8+ has a prior of $1.60\%$ and forces mate. Here is what PUCT does with it, from the actual
engine:

\begin{realdata}[title={Real engine output: Morphy position, ranked by PUCT score $S$}]
At $N(s) = 182$ visits (search still running):
\begin{center}
\begin{tabular}{@{}lrrrrr@{}}
\toprule
Move & $P$ & $N$ & $Q$ & $U$ & $S = Q+U$ \\
\midrule
Qb7   & 31.93\% & 169 & $0.6575$ & $0.0445$ & $\mathbf{0.70199}$ \\
\textbf{Qb8+} & \textbf{1.60\%} & \textbf{0} & $0.3187^{*}$ & $0.3791$ & $\mathbf{0.69780}$ \\
Rxd7  &  7.81\% &   4 & $0.3155$ & $0.3701$ & $0.68387$ \\
Qc3   &  8.33\% &   2 & $-0.0005$ & $0.6583$ & $0.65774$ \\
Qb5   &  9.57\% &   2 & $-0.1019$ & $0.7565$ & $0.65235$ \\
\bottomrule
\end{tabular}
\end{center}
\footnotesize $^{*}$ the FPU value recovered in §6.4.2; the engine's \texttt{Q} column prints
$0.0$ for unvisited moves.
\normalsize

\medskip
The mating move, which the network rated at $1.60\%$ and which has \textbf{never been
looked at}, is ranked \emph{second} --- behind the leader by $0.004$. Its entire claim comes
from the $U$ term: $169$ visits have piled into Qb7, and $\sqrt{N}/(1+n)$ has been quietly
compounding.

Six visits later, the search selects it. The result:
\begin{center}
\begin{tabular}{@{}lrrrrr@{}}
\toprule
At $N(s) = 188$ & $P$ & $N$ & $Q$ & $U$ & $S$ \\
\midrule
\textbf{Qb8+} & 1.60\% & \textbf{3} & $\mathbf{1.0000}$ & $0.0964$ & $\mathbf{1.09636}$ \\
Qb7   & 31.93\% & 172 & $0.6587$ & $0.0445$ & $0.70310$ \\
\bottomrule
\end{tabular}
\end{center}
$Q = 1.0$ exactly: not an estimate but a proof --- 16.Qb8+ Nxb8 17.Rd8\# is a terminal
sequence, so the value is certain rather than sampled. The engine reports
\texttt{bestmove b3b8} and \emph{stops searching}, at 188 nodes out of a requested 6400.
\end{realdata}

\begin{keyidea}
Three visits beat one hundred and seventy-two. This is worth dwelling on, because it is where
the naive reading of Monte-Carlo search --- "the most-visited move wins" --- breaks. A proven
result outranks a well-sampled estimate. Chapter~\ref{ch:readingtree} deals with the general
question of which number decides the move, and Chapter~\ref{ch:engineering} with how
certainty propagates.
\end{keyidea}

And notice what would have happened under each rejected attempt: Attempt~1 (top-3 filter)
never sees the move. Attempt~2 (constant prior bonus) gives it a standing $+0.016$ against
Qb7's $+0.319$ --- it would need the search to have already discovered the truth it is
supposed to be discovering. Attempt~4 ($\sqrt{\ln N}$ numerator) reaches $U \approx 0.04$
instead of $0.379$ at the same point, ranking Qb8+ somewhere in the twenties. \textbf{The
$\sqrt{N}$ is the reason the mate was found.}

\begin{pitfall}
Do not over-generalise this happy ending. PUCT guarantees that a low-prior move's claim
\emph{grows}, not that it grows \emph{in time}. If the prior had been $0.05\%$ rather than
$1.60\%$, the same computation gives $U \approx 0.012$ at $N=182$, and the move would need
roughly a thousand times more search to surface. Leela does miss things, and the shape of
what it misses is precisely "moves its policy rates near zero in positions where nothing else
is failing". Chapter~\ref{ch:lookahead} is about whether we can detect those cases from
inside the network instead of waiting for search to stumble on them.
\end{pitfall}

\begin{repobox}[title={in this repository}]
$P$ is the \texttt{policy} array this project already extracts and draws as arrows on the
board; \texttt{U} and \texttt{S} are visible in \texttt{--verbose-move-stats}. The
"policy-blindness" metric in your diagnosis layer is a statement about $P$; the
"optical trap / hidden gem" classification is a statement about the gap between $P$ and the
visit distribution, i.e.\ about how much of §6.3's sharpening happened and in which direction.
\end{repobox}

\begin{exercises}

\exhead[what P means]{ch6:pmeans} In one sentence each, state what $P(a\given s) = 0.32$
claims and what it does \emph{not} claim. Then explain how a move can have $P = 1.6\%$ and be
the only move that wins.

\exhead[compute U]{ch6:computeu} A node has $N(s) = 100$ total visits. Move $a$ has
$P = 0.25$ and $N(s,a) = 40$; move $b$ has $P = 0.05$ and $N(s,b) = 0$. Using
$c_{\text{puct}} = 1.745$, compute $U$ for both. Which is larger, and by what factor?

\exhead[the crossover]{ch6:crossover} Continue the previous exercise: move $a$ has
$Q = 0.55$ and move $b$ has $Q_{\text{FPU}} = 0.30$. Compute $S$ for both. Now recompute at
$N(s) = 400$ with $a$ having absorbed all the extra visits ($N(s,a) = 340$). At which of the
two search sizes does $b$ get selected? State the mechanism in one sentence.

\exhead[verify the engine]{ch6:verify} From the real table in §6.5, take the Qb7 line at
$N(s)=182$: $P = 0.3193$, $N(s,a) = 169$. Compute $U$ using Equation~\eqref{eq:puct} with
$c_{\text{puct}}$ from Equation~\eqref{eq:cpuctgrow}. Compare with the engine's $0.04453$.
How many decimal places do you match?

\exhead[the equilibrium theorem]{ch6:equilibrium} Reproduce the proof of §6.3 in your own
words. Then: a node has three moves with priors $0.6, 0.3, 0.1$ and all $Q$ equal. After 100
visits, what does the theorem predict for $n_1, n_2, n_3$? (Remember it is $1+n_a$ that is
proportional to $P_a$; discuss whether the "+1" matters at 100 visits and at 3 visits.)

\exhead[sharpening]{ch6:sharpening} Using the real 1600-node table for the initial position,
compute for each move the ratio (visit share)/(prior). Which moves have ratio above 1?
Explain what a ratio of $1.60$ for e4 and $0.43$ for e3 tells you, in terms of §6.3.

\exhead[FPU by hand]{ch6:fpu} A node has $Q(s) = -0.20$ and has visited moves covering $45\%$
of the prior mass. Compute $Q_{\text{FPU}}$ for its unvisited moves. Now suppose the visited
mass rises to $90\%$: recompute. Explain in one sentence why the penalty should grow.

\exhead[FPU's effect on style]{ch6:fpustyle} Suppose you set \texttt{FpuValue} to $0.0$
(unvisited moves inherit the parent's value untouched). Describe how the search's behaviour
changes: broader or narrower? Now set it to $1.0$. Which setting would make an engine more
likely to find a surprising sacrifice, and what does it cost?

\exhead[$c_{\text{puct}}$ growth]{ch6:cpuct} Using Equation~\eqref{eq:cpuctgrow} with your
engine's defaults, compute $c_{\text{puct}}$ at $N = 800$, $N = 10^5$, and $N = 10^7$. Plot
roughly. What does the designers' choice say about how much to trust the policy at each
scale, and does it agree with the answer you gave in Exercise~\ref{ex:ch4:tuningc}?

\exhead[the hidden move's timetable]{ch6:timetable} A move has $P = 0.005$ and $N(s,a)=0$ in
a node where the leader holds $Q = 0.60$ and the FPU value is $0.30$. Using
$c_{\text{puct}} \approx 1.75$, solve for the value of $N(s)$ at which the unvisited move's
$S$ overtakes the leader's $Q$. Comment on whether that node count is reachable in your
tool's configuration, and what this implies for what the tool can promise a user.

\exhead[design]{ch6:design} Compare Equation~\eqref{eq:puct} with the formula you invented in
Exercise~\ref{ex:ch5:designfix}. Name one thing yours does better and one thing it does
worse. (This is not a trick question --- several published variants differ from PUCT in
defensible ways.)

\end{exercises}
